% !TeX root = ../papers/nc-p3.tex
\section{有限窓の頂点次元二の単純と二次曲面の伝播}
\label{sec:finite-two-dimensional-quadrics}
\begin{theorem}\label{thm:finite-two-dimensional-quadrics}
全四辺が有限な標準有限Frobenius窓の、各頂点次元2の折り畳み単純加群を考える。
同型と共通次数倍率を除いた分類は、各文字空間 $V_i$ 上の階数3または4の射影二次式 $q_i$ と、
階数4のときの一方の罫線族の指定による次の有限方程式の解に一致する。
二行二列線型表示 $A_i:V_i\to\mathrm{Hom}(U_i,W_i)$、$\det A_i=q_i$ をその指定に従って一つ取る。
各辺で $C_i:W_i\to U_{i+1}$ を未知とする線型写像
\[
 C_i\longmapsto\left(\sum_{u,v} r_{uv} A_{i+1}(v) C_i A_i(u)\right)_{r\in R_i}
\]
が非零核を持つことが条件である。核があれば次元1で可逆行列からなり（零元を除く）、
四箇所の間に追加の接続条件はない。この判定は未認証の窓と非被約な点グラフにも適用できる。
\end{theorem}
\begin{proof}
まず任意の単純で全頂点が非零なら、四矢に共通する核も、全像の張る真部分空間もない。
前者は頂点単純部分、後者は頂点単純商を与えるためである。
あるベクトルの次位置での像が直線なら矢の作用はその上で非零共変元 $p_i$ を通る。
六関係をこのベクトルに作用させると、次の直線は $R_i(p_i,-)$ で消える。
有限辺でこの評価の階数は少なくとも3なので、消す三平面が定まり、次の像も直線であり、
隣接点は有限グラフの一意な後続 $p_{i+1}$ になる。
四移動はスキーム上で恒等である。各位置で $\ker p_i$ の全矢に消されるベクトル空間を $K_i$ とすると、
$p_i$ で値1の矢は $K_i\to K_{i+1}$ に単射である。共通核の不在と関係評価から従う。
従って四段で同型となり、合成の固有直線を選べば四頂点一次元の部分加群を得る。
全頂点次元2の単純では矛盾する。ゆえに各 $0\ne u\in U_i$ に対し $A_i(V_i)u=W_i$ である。
逆にこの条件が全辺で成立すれば、任意の非零部分加群は次位置を全て含み、巡回して全体に一致する。

二行二列行列空間でこの条件は、像の直交空間が階数1の元を含まないことと同値である
（$w^*A_i(V_i)u=0$ を考える）。行列式二次式のため、直交空間は零か可逆元の張る直線だけである。
後者の像は左右基底変更で跡零超平面にでき、行列式はそこに非退化に制限される。
従って $\det A_i$ の階数はそれぞれ4または3である。
階数4では滑らかな二次曲面の二つの罫線を区別したSegre表示が、左右二次元空間の基底変更まで
行列表示を一意に復元する。階数3では根基を割って滑らかな平面二次曲線を取り、跡零三空間と同定する。
その二次曲線の全射影自己同型は二次元空間の共役作用で誘導されるので表示は一通りである。
射影倍率はいずれも左右変更に吸収できる。

接続方程式の核に階数1の $a b^*$ があれば、$u\in U_i$ と $e^*\in W_{i+1}^*$ を動かした
$b^* A_i(-)u, e^* A_{i+1}(-)a$ の全組が辺の点対応に入る。
両方とも二次元の変数から文字双対への単射である。
従って直線二本の積を含み、有限性に矛盾する。
核の二平面は行列式二次式の零点を持つので、可逆な解直線が唯一となる。
四接続を選べば関係は全て満たされ、上の像条件で単純である。
接続の四倍率は頂点基底の倍率で積だけになり、体が代数閉なので残りは次数倍率で吸収できる。
\end{proof}

\begin{proposition}\label{prop:two-dimensional-quadric-successors}
一つの向き付き $q_i$ だけを与えて次の二次曲面を列挙できる。
$A=A_i$ とし、
\[
 K=\ker\left[V_{i+1}^*\otimes W_i^*\longrightarrow R_i^*\otimes U_i^*\right],
 \qquad b\longmapsto\left(\sum r_{uv} b(v)A(u)\right)_r .
\]
$\dim K\le3$ であり、次位置が存在するのは $\dim K=2,3$ のときだけである。
次位置は $K$ の二平面を選び、それを行とする行列の行列式と罫線指定を取ったものに限り、全てがこの一辺を満たす。
次元2なら一意である。次元3なら次の二次曲面は、$\mathbb P(V_{i+1})$ 内の一つの滑らかな空間三次曲線を
通る二次曲面の網全体を一度ずつ動く。この手順を四辺で閉じることも単純の分類となる。
\end{proposition}
\begin{proof}
$K$ に非零可分解元 $p\otimes b^*$ があれば、前証明と同じく
$(b^*A(-)u,p)$ の直線が有限辺に入ってしまう。
$\mathbb P(W_i^*\otimes V_{i+1}^*)$ のSegreは余次元3なので $\dim K\le3$。
選んだ二行の全非零線型結合は行列階数2である。従ってこれらの行が作る次の矢の像は
全非零ベクトル上で全射で、前定理の階数3または4の二次曲面になる。
逆に後続の行は接続で $W_i$ 上に移せば必ずこの形である。

三行全部を束ねた $B:V_{i+1}\to\mathrm{Hom}(W_i,K^*)$ を考える。
どの $0\ne w\in W_i$ でも $B(-)w$ は三次元へ全射で、核は四次元文字空間の一直線である。
$B$ が単射でなければその共通核を除く三次元からの全射となり、$w$ に線型な三行三列の行列式が
射影直線上で消えないことになり矛盾する。よって核直線の射 $\mathbb P(W_i)\to\mathbb P(V_{i+1})$
の異なる二点は異なる像を持つ。微分も単射である。零なら一次変分でも同じ文字が
独立な二つの $w$ を消し、$B$ の単射性に反する。三次小行列は基点のない三次表示を与え、
従って像は滑らかな空間三次曲線である。$B(v)$ の三つの二次小行列はこの曲線のイデアルを生成する
（階数一条件が曲線をスキームとして定めることは各非零成分で核を解いてもよく、あるいは余次元2の
Hilbert--Burch表示を用いる）。これらは三独立で、二行の外積が三小行列の全一次結合を一意に与える。
\end{proof}
